%% Introduction (MedIA revision, v2 pipeline). Numbers follow Section 5 (sections/results.tex).
\section{Introduction}
\label{sec:intro}

Deep networks for histopathology often lose accuracy when they are applied to
slides from a laboratory, scanner or patient population they were not trained
on. Staining protocols, scanners and tissue handling differ between sites, and
networks pick these differences up: the submitting site of a TCGA slide can be
predicted from the image, and this signal biases downstream predictions
\citep{howard2021site}. Color differences between laboratories degrade
classifiers on unseen laboratories \citep{tellez2019stain}, and mitotic-figure
detectors deteriorate on new scanners and tumor types
\citep{aubreville2023midog,aubreville2024midog2022}. The problem is not
specific to pathology; most externally validated radiology models also perform
worse on external data \citep{yu2022external}.

For practitioners, the goal is a model that is trained once and then works at
every hospital where it is used, without collecting labels or retraining at
each new site. Many remedies promise this: stain normalization and color
augmentation, domain-generalization (DG) objectives, self-supervised
pre-training, test-time adaptation and, most recently, pathology foundation
models. To choose between them, practitioners need evidence that can tell a
real difference from noise.
That is harder than it looks. Benchmarks with a single official test site,
such as Camelyon17-WILDS, evaluate every method on one domain, and ERM accuracy
on that domain has a standard deviation of 6.4 points across ten seeds
\citep{koh2021wilds}.
Across DG benchmarks, ERM is hard to beat once every method is tuned and
evaluated under the same protocol \citep{gulrajani2021domainbed}. Rankings in
biomedical image analysis challenges are sensitive to the test data and the
ranking scheme \citep{maierhein2018rankings}. A reported gain can thus reflect
the chosen test site, the seed or the tuning budget as much as the method.

This study began with such a gain. In an earlier version of this work, a
single Camelyon17 test hospital suggested that self-supervised pre-training and
sharpness-aware minimization were synergistic. Each method alone was worse than
ERM, and together they were better. Repeating the runs and holding out every
hospital in turn reversed the sign of this interaction at most hospitals
(Section~\ref{sec:casestudy}). The same earlier version argued that the type of
shift, scanner-only or institutional, decides which method works best. That
argument rested on two cohorts that differed in task as well as in shift type,
and an audit of our code found defects in several method implementations. We
therefore rebuilt the study around six questions:
\begin{enumerate}
  \item How much of the variance of a held-out score comes from the choice of
    held-out domain, and how much from run-to-run noise?
  \item With the few domains that pathology cohorts provide, can realistic
    differences between methods be detected at all?
  \item How stable are method rankings within a cohort?
  \item Do rankings carry over between cohorts that share the task, or between
    cohorts that share the type of shift?
  \item Which choices give a model the best chance of working at an unseen site
    without retraining?
  \item How can a practitioner decide, with few local labels, whether a trained
    model is good enough at a particular new site?
\end{enumerate}

To answer them, we use a \emph{domain-complete, seed-replicated} protocol.
Every domain of a cohort is the test domain of one leave-one-domain-out fold,
every method is run at several seeds under one fixed training budget, and
hyperparameters and checkpoints are selected without test data. We analyze the
results with a random-effects decomposition into between-domain and
between-run variance with interval estimates, paired contrasts with a
fold-dependence sensitivity analysis, a power calculation based on the
variance of the paired difference, and bootstrap intervals for ranks and for
rank agreement between cohorts. We apply the protocol to four public cohorts
arranged in a $2\times2$ design. Two tumor-classification cohorts and two
mitotic-figure cohorts each pair an acquisition-only shift (scanners within
one laboratory) with an institutional shift (hospitals or tumor types).
We evaluate up to 21 arms at two model scales, a compact CNN and ImageNet
ResNet-50, together with frozen and fine-tuned pathology foundation models,
in 4,778 training runs. No model is ever trained on labeled data from its test
site, and apart from the test-time adaptation arms, which use unlabeled test
images, every result measures how a model trained once performs, unchanged, at
a site it has never seen.

Our contributions are:
\begin{itemize}
  \item an empirical-Bayes \emph{site acceptance test} that decides whether a
    trained model reaches a target AUROC at a new site. It combines the
    between-site spread of the development results with a small labeled sample
    from the new site, plans the number of labels in advance, and is validated
    on every held-out site of the study (Sections~\ref{sec:sat}
    and~\ref{sec:results-sat});
  \item a six-step, evidence-graded roadmap for practitioners who need models
    that work across hospitals and scanners without retraining, covering the
    target definition, data collection, training recipe, label-free adaptation
    at a new site, validation and site acceptance
    (Section~\ref{sec:roadmap});
  \item an evaluation protocol and a set of diagnostics for distribution shift
    in multi-domain pathology cohorts, with the code, the fold definitions and
    every per-run record released so that each table and figure can be
    regenerated;
  \item a $2\times2$ task-by-shift study at two model scales, with
    domain-generalization objectives, stain and color methods,
    self-supervision, test-time adaptation, pathology foundation models and
    two controls for self-supervision, all tuned per fold;
  \item an external reference point for the variance decomposition, computed
    from DomainBed's published per-environment results;
  \item two negative results reported in full: a single-hospital synergy
    claim that does not survive the protocol, and a single-institution cohort
    that cannot support domain claims.
\end{itemize}

The main findings are these. Without retraining, the best standard-scale
method of each cohort reached a worst-site AUROC of 0.884 to 0.986, against
0.534 to 0.895 for ERM, but no method was best everywhere. For most arms, the held-out domain accounts for
most of the variance of a single run (median ICC 0.72--0.95), and the worst
domain differs between methods. After correction for multiple comparisons, no
standard-tier method beats ERM on any cohort. Even a mean gain of 0.145 AUROC
that holds at every domain is not significant with five domains, and detecting
a gain of 0.05 on the tumor cohorts would need between 7 and 53 domains. Under
the primary selection rule, none of the six DG objectives ranks in the top five
on any cohort. At the standard
tier, method rankings agree more between cohorts that share a task than
between cohorts that share a shift type. With one cohort per cell this is a
description, not a proof, and the compact tier does not clearly show it. Fine-tuned
pathology foundation models with color augmentation were the most reliable
choice on average, and re-estimating batch-normalization statistics at the new
site, which needs no labels and no weight updates, helped mainly under severe
scanner shift. Implementation choices such as per-tile versus per-slide
stain estimation, and the learning rate of a pretrained encoder, change
results as much as the choice of method. Finally, the development results
alone would have accepted 8.6\% of the sites that fell below an AUROC of 0.80.
At standard scale, the site acceptance test kept false acceptance at or below
2.1\% and, for targets of 0.80 and 0.85, reached the same power as local labels
alone with about half as many labels.

Section~\ref{sec:related} reviews related work. Section~\ref{sec:protocol}
defines the protocol, the statistics and the site acceptance test, and
Section~\ref{sec:setup} the cohorts, models and arms. Section~\ref{sec:results} reports the results.
Section~\ref{sec:discussion} discusses them, presents the roadmap
(Section~\ref{sec:roadmap}) and the two case studies, and lists the
limitations.
